As chiplet-based design has moved from a promising concept toward a practical system-integration strategy, the communication fabric has become a central architectural concern. Open die-to-die standards such as UCIe now define not only physical signaling, but also protocol, compliance, and system-level interoperability for multi-chip systems, while recent chiplet research emphasizes that high bandwidth, low latency, and efficient coordination across heterogeneous dies are now first-order design requirements \cite{9773184}. These requirements are especially difficult to satisfy in heterogeneous systems, where CPUs, memories, and accelerators generate traffic with different quality-of-service demands; earlier NoC research already showed that conventional on-chip interconnects may not serve latency-sensitive CPU traffic and bandwidth-intensive accelerator traffic equally well.As a result, current work increasingly explores topology-aware and application-adaptive communication fabrics, including dynamically tailored sub-networks, optimized chiplet arrangements, and flexible die-to-die interfaces that better match communication structure to workload behavior \cite{chen2024surveychipletbasedintegratedarchitecture}. 

At the same time, the field is moving in two partially competing directions. On one hand, standardized chiplet interfaces promise portability, vendor interoperability, and ecosystem-level reuse; on the other, their expanding support for test, debug, manageability, and packaging scalability implies a broader implementation scope that is not always necessary for smaller research platforms \cite{9773184}. For this reason, lightweight NoC-based communication remains technically relevant: a compact mesh router can expose the essential problems of routing, arbitration, buffering, and flow control without the full complexity of a production die-to-die stack. FPGA prototyping is particularly valuable in this context, because both research on NoC emulation and current hardware-assisted verification practice highlight its importance for faster hardware evaluation, system validation, and performance observation beyond software-only simulation \cite{10035152}.

This thesis follows that implementation-oriented perspective by developing a hardware NoC router for a $3 \times 3$ mesh and realizing it as an FPGA-oriented prototype. The design supports deterministic XY routing, virtual-channel flow control, pipelined forwarding, arbitration among competing packet flows, and a defined flit format for packet transfer. By integrating multiple router instances into a nine-node mesh, the work provides a concrete platform for observing multi-hop communication, contention, and flow-controlled packet delivery in hardware rather than only through abstract architectural modeling.
